\section*{Bab \ref{ch:g7:identifying-substances} --- Mengenali Zat}

\begin{solution}{exo:g7:identifying-substances:1}
Uji air kapur: gas itu dialirkan melalui air kapur yang jernih, yang
menjadi keruh jika gas itu karbon dioksida.
\end{solution}

\begin{solution}{exo:g7:identifying-substances:2}
Serbuk itu tembaga sulfat anhidrat, dan cairan itu mengandung air.
\end{solution}

\begin{solution}{exo:g7:identifying-substances:3}
Oksigen.
\end{solution}

\begin{solution}{exo:g7:identifying-substances:4}
Pensil tidak \omterm{def:g2:mixing-and-dissolving:dissolve}{larut} dalam \omterm{def:g6:solutions-and-solubility:solute}{pelarut}. Garis tinta sendiri akan terbawa naik
dan terpisah, sehingga merusak \omterm{def:g7:identifying-substances:chromatography}{kromatogramnya}.
\end{solution}

\begin{solution}{exo:g7:identifying-substances:5}
Jika totolannya berada di bawah \omterm{def:g6:solutions-and-solubility:solute}{pelarut}, zat-zat warnanya akan \omterm{def:g2:mixing-and-dissolving:dissolve}{larut} ke
dalam \omterm{def:g6:solutions-and-solubility:solute}{pelarut} di gelas kimia alih-alih naik melalui kertas.
\end{solution}

\begin{solution}{exo:g7:identifying-substances:6}
Bukan, tinta itu \omterm{def:g6:pure-substances-and-mixtures:mixture}{campuran} dua zat warna: kemungkinan zat warna kuning
dan zat warna biru pembanding.
\end{solution}

\begin{solution}{exo:g7:identifying-substances:7}
Tampung sedikit gas itu dalam tabung reaksi dan dekatkan bilah kayu
menyala ke mulutnya: letupan melengking menunjukkan hidrogen.
\end{solution}

\begin{solution}{exo:g7:identifying-substances:8}
\omterm{def:g4:air-a-mixture-of-gases:air}{Udara} yang diembuskan mengandung jauh lebih banyak karbon dioksida
daripada \omterm{def:g4:air-a-mixture-of-gases:air}{udara} segar.
\end{solution}

\begin{solution}{exo:g7:identifying-substances:9}
Uji itu menunjukkan bahwa air ada, bukan bahwa air itu sendirian: air
asin, jus, atau \omterm{def:g6:pure-substances-and-mixtures:mixture}{campuran} berair apa pun juga akan membirukan serbuk itu.
\end{solution}

\begin{solution}{exo:g7:identifying-substances:10}
Tiga zat warna (tiga bercak). Zat warna kuning naik paling tinggi.
\end{solution}

\begin{solution}{exo:g7:identifying-substances:11}
Misalnya: bilah kayu berbara dimasukkan ke setiap tabung --- yang
menyalakannya kembali berisi oksigen. Pada dua tabung lainnya, bilah
kayu menyala didekatkan ke mulutnya --- letupan melengking menunjukkan
hidrogen. Tabung terakhir berisi karbon dioksida (dapat dipastikan
dengan air kapur). Dua atau tiga uji sudah cukup.
\end{solution}

\begin{solution}{exo:g7:identifying-substances:12}
Gas itu mungkin tidak mengandung karbon dioksida, atau terlalu sedikit
untuk terlihat dalam waktu itu. Uji kontrol dengan \omterm{def:g4:air-a-mixture-of-gases:air}{udara} embusan napas,
yang diketahui mengandung karbon dioksida, memeriksa bahwa air kapurnya
berfungsi.
\end{solution}

\begin{solution}{pb:g7:identifying-substances:1}
\textbf{1.} Agar kondisinya (kertas, \omterm{def:g6:solutions-and-solubility:solute}{pelarut}, waktu) sama untuk
keempatnya: hanya dengan begitu ketinggian bercak-bercaknya dapat
dibandingkan.

\textbf{2.} \omterm{def:g6:pure-substances-and-mixtures:mixture}{Campuran}: tintanya memberi tiga bercak terpisah, tiga zat
warna.

\textbf{3.} Dua zat warna. Tidak: tinta catatan itu memiliki tiga zat
warna, dan pena 1 tidak memiliki zat warna merah maupun kuning.

\textbf{4.} Ketiga bercaknya tidak setinggi bercak-bercak catatan itu:
ketiganya zat warna yang berbeda.

\textbf{5.} Pena 2: ketiga bercaknya tepat setinggi ketiga bercak
catatan itu.

\textbf{6.} Oksigen.

\textbf{7.} Hidrogen.

\textbf{8.} Karbon dioksida.

\textbf{9.} Blanko (sebuah kontrol). Uji ini menunjukkan bahwa air kapur
tidak menjadi keruh dengan sendirinya, ataupun dengan \omterm{def:g4:air-a-mixture-of-gases:air}{udara} biasa dalam
waktu itu: jadi kekeruhan dengan gas Z memang berasal dari karbon
dioksida.

\textbf{10.} ``Air --- tembaga sulfat anhidrat --- putih menjadi biru'':
cairan itu mengandung air.

\textbf{11.} Tidak. Uji yang positif menunjukkan bahwa air ada, bukan
bahwa tidak ada yang lain: cairan itu bisa saja larutan atau \omterm{def:g6:pure-substances-and-mixtures:mixture}{campuran}
berair yang lain.

\textbf{12.} Tinta catatan itu mengandung 3 zat warna, dan pena 2 yang
menulis catatan itu.
\end{solution}
